\section{GAN：生成对抗网络}

\subsection{从对抗攻击到生成模型}

GAN 的思想起源于对抗攻击（Adversarial Attack）。授课教师通过一个经典例子说明：一个训练好的分类器可以正确识别熊猫，但如果在图像上添加一个人眼几乎不可见的微小扰动（噪声），分类器就会将其误判为鸵鸟——即使将差异放大 10 倍，人眼也很难看出区别。

\begin{knowledgebox}{从对抗攻击到 GAN 的思想飞跃}
对抗攻击揭示了一个事实：神经网络可以对人眼无法察觉的细微变化极其敏感。GAN 的创始人 Ian Goodfellow 正是受此启发：如果一个网络（攻击者/生成器）可以生成``欺骗''另一个网络（防御者/判别器）的样本，那么通过两者的竞争，生成器最终能学会生成逼真的数据。
\end{knowledgebox}

\subsection{GAN 的架构}

GAN 由两个神经网络组成：

\begin{itemize}
    \item \textbf{Generator（生成器）} $G_\theta$：接受 latent 样本 $z \sim p(z)$，输出``假''数据点 $G_\theta(z)$
    \item \textbf{Discriminator（判别器）} $D_\phi$：接受一个数据点，输出该数据点为``真实数据''的概率 $D_\phi(x) \in {[}0, 1{]}$
\end{itemize}

生成器和判别器互相竞争：
\begin{itemize}
    \item 判别器的目标：正确区分真实数据和生成数据
    \item 生成器的目标：生成能骗过判别器的样本
\end{itemize}

\subsection{GAN 的损失函数：Minimax 博弈}

GAN 的训练目标是一个 minimax 优化问题：

$$\min_\theta \max_\phi \; V(G_\theta, D_\phi) = \mathbb{E}_{x \sim p_{\text{data}}} {[}\log D_\phi(x){]} + \mathbb{E}_{z \sim p(z)} {[}\log(1 - D_\phi(G_\theta(z))){]}$$

\begin{importantbox}{GAN 损失函数各项的含义}
\begin{itemize}
    \item $\mathbb{E}_{x \sim p_{\text{data}}} {[}\log D_\phi(x){]}$：对真实数据，判别器应输出高概率（接近 1），使 $\log D_\phi(x)$ 尽量大
    \item $\mathbb{E}_{z \sim p(z)} {[}\log(1 - D_\phi(G_\theta(z))){]}$：对生成数据 $G_\theta(z)$，判别器应输出低概率（接近 0），使 $1 - D_\phi(G_\theta(z))$ 接近 1
\end{itemize}
\textbf{判别器}通过\textbf{最大化} $V$ 来提高区分能力；\textbf{生成器}通过\textbf{最小化} $V$ 来提高欺骗能力。
\end{importantbox}

\subsubsection{判别器的最优解}

固定生成器 $G_\theta$，对判别器求最优解：

$$D^*_\phi(x) = \frac{p_{\text{data}}(x)}{p_{\text{data}}(x) + p_G(x)}$$

其中 $p_G(x)$ 是生成器诱导的分布。当 $p_G = p_{\text{data}}$ 时，$D^*(x) = 1/2$，即判别器无法区分真假——这正是 GAN 训练收敛的理想状态。

\subsubsection{全局最优的 GAN}

可以证明，当 $p_G = p_{\text{data}}$ 时，损失函数达到全局最小值 $-2\log 2$。此时生成器的分布与数据分布完全匹配。

\subsection{GAN 的训练难题}

\begin{warningbox}{GAN 训练中的常见问题}
Minimax 优化问题在实践中\textbf{极其难以优化}，容易陷入糟糕的局部最优。常见的失败模式包括：
\begin{itemize}
    \item \textbf{判别器过强（Discriminator Dominance）}：判别器收敛速度远快于生成器，导致生成器得到的梯度信号极弱（几乎为零），无法有效学习。这是最常见的失败模式。
    \item \textbf{Mode Collapse（模式坍塌）}：即使训练顺利进行，生成器也可能只学会生成数据分布中的部分模式（modes），丢失了数据的多样性。例如在多模态数据上，生成器可能只集中生成某一个模式的样本。
    \item \textbf{训练不稳定}：loss 曲线剧烈震荡，难以判断是否收敛。
\end{itemize}
\end{warningbox}

授课教师提到了一个生动的类比：如果你一开始就设定了极其严格的成功标准（discriminator 太强），那就什么也学不到了——这就像生活中的道理一样。

\subsubsection{GAN 训练的工程实践}

在实践中，GAN 训练需要大量的工程技巧：

\begin{itemize}
    \item 通常先预训练生成器几步，让它能产出略有意义的输出，再引入判别器进行对抗训练
    \item 需要仔细平衡两个网络的学习率和训练步数比例
    \item 许多公司在 2020 年之前投入了大量专门的工程团队来调试 GAN 训练参数
\end{itemize}

\subsection{GAN 的发展历程}

GAN 自 2013/2014 年由 Goodfellow 等人提出后，涌现了大量变体。授课教师简要提到了几个里程碑式的工作：

\begin{table}[H]
\centering
\begin{tabular}{lll}
\toprule
模型 & 年份 & 主要贡献 \\
\midrule
GAN & 2014 & 提出对抗训练框架 \\
DCGAN & 2015 & 引入卷积架构 \\
WGAN & 2017 & 用 Wasserstein 距离替代 JS 散度 \\
Progressive GAN & 2017 & 渐进式增长训练策略 \\
StyleGAN / StyleGAN2 & 2019/2020 & 高质量人脸生成 \\
\bottomrule
\end{tabular}
\caption{GAN 的重要变体（部分）}
\end{table}

\begin{knowledgebox}{为什么从 GAN 转向 Diffusion Model？}
授课教师指出，业界从 GAN 转向 Diffusion/Flow Model 的主要原因并非质量——StyleGAN 在 2019 年已经能生成非常逼真的人脸图像。真正的驱动力是：（1）\textbf{训练困难}——minimax 优化极不稳定，需要大量工程投入；（2）\textbf{多样性不足}——mode collapse 问题导致生成样本缺乏多样性。Diffusion Model 通过将生成过程分解为多步去噪，巧妙地规避了这两个问题。
\end{knowledgebox}

\subsection{本章小结}

GAN 通过生成器和判别器的对抗博弈来学习数据分布，其理论框架优美（minimax 博弈在理想情况下能精确匹配数据分布），但实践中面临训练不稳定、mode collapse 等严重问题。这些问题促使了 VAE 和后来的 Diffusion Model 等替代方案的发展。

\newpage
